\documentclass[11pt]{amsart}

\usepackage[T1]{fontenc}
\usepackage{lmodern}
\usepackage{microtype}
\usepackage{mathtools,amssymb,amsthm}
\usepackage[hidelinks]{hyperref}

\hypersetup{
  pdftitle={A 71-point arc refutes the printed finite-p form of the entropic
            Cauchy-Davenport candidate of Gavalakis, Goh and Kontoyiannis}
}

\newtheorem{theorem}{Theorem}
\newtheorem{lemma}[theorem]{Lemma}
\newtheorem{proposition}[theorem]{Proposition}
\newtheorem{corollary}[theorem]{Corollary}
\theoremstyle{remark}
\newtheorem*{remark}{Remark}

\DeclareMathOperator{\Ent}{H}
\newcommand{\Fp}{\mathbb F_p}

\title[A 71-point arc refutes the printed finite-$p$ form of a candidate]
{A 71-Point Arc Refutes the Printed Finite-$p$ Form of the Entropic
 Cauchy--Davenport Candidate of Gavalakis, Goh and Kontoyiannis}

\date{}

\subjclass[2020]{94A17, 11B13, 11B30, 05B10}
\keywords{Shannon entropy, Cauchy--Davenport theorem, sumset,
entropy power inequality, prime field, counterexample}

\begin{document}

\begin{abstract}
Gavalakis, Goh and Kontoyiannis suggest, in a hedged display whose first line
is the only part at issue here, that a candidate entropic Cauchy--Davenport
theorem ``might be'' that
\[
  \Ent(X+X')\ \ge\ \min\bigl\{\log p,\
  \log\bigl((2^{\Ent(X)+1}-1)/\sqrt2\bigr)\bigr\}
\]
for every prime $p$, every $\Fp$-valued $X$ and $X'$ an independent copy of
$X$, the logarithm being to base~$2$.  We refute that first line read literally at
finite $p$, in the branch where the minimum is \emph{not} the cap $\log p$:
take $p=101$ and
$X$ uniform on the arc $S=\{0,1,\dots,70\}$.  Then $2^{\Ent(X)+1}-1=141$
exactly, the integer comparison $141^{2}=19881<20402=2\cdot101^{2}$ shows that
the active branch is $\log(141/\sqrt2)$ and not $\log 101$, and
$\Ent(X+X')<\log(141/\sqrt2)$.  The decisive step is a single comparison
between two explicit positive integers; no logarithm is evaluated and no
floating-point arithmetic is used.  Six further witnesses
at $p=61,151,211,307,401,509$ are certified the same way.  Only the finite-$p$
first line of the display is affected: the asymptotic second line, which
carries an $o_{p\to\infty}(1)$, is untouched, and nothing here contradicts the
authors' own theorem, an $\varepsilon$-weakened statement restricted to an
entropy window: we do not determine that window's constant, and what we check
is that all seven witnesses satisfy the theorem's conclusion for every
$\varepsilon>0.023$.
\end{abstract}

\maketitle

\section{The statement}

Throughout, $\log$ is the logarithm to base~$2$, entropies are in bits, and
$0\log 0=0$, as in the source paper \cite{GGK}.  For a prime $p$ we write
$\Fp$ for the field with $p$ elements, and $X'$ always denotes an independent
copy of the $\Fp$-valued random variable $X$.

Gavalakis, Goh and Kontoyiannis \cite{GGK} discuss what an entropic analogue
of the Cauchy--Davenport theorem $|A+B|\ge\min\{p,|A|+|B|-1\}$ could say.
Having observed that the naive analogue fails, they write ``In view of the
above comments, a more appropriate candidate for an entropic
Cauchy--Davenport theorem might be that'', and then display, as
equation~(2) of \cite{GGK},
\begin{align*}
  \Ent(X+X')\ &\ge\ \min\left\{\log p,\
   \log\!\left(\frac{2^{\Ent(X)+1}-1}{\sqrt2}\right)\right\}\\
  &=\ \min\left\{\log p,\ \Ent(X)+\tfrac12-o_{p\to\infty}(1)\right\}.
\end{align*}
The first line of that display, which is the only part at issue here, reads
\begin{equation}\label{eq:cand}
  \Ent(X+X')\ \ge\ \min\left\{\log p,\
  \log\!\left(\frac{2^{\Ent(X)+1}-1}{\sqrt2}\right)\right\}.
\end{equation}
Read literally, the quantifiers are the ones the display carries and no
others: every prime $p$, every $\Fp$-valued $X$, and $X'$ an independent
copy; there is no largeness hypothesis on $p$, no $\varepsilon$, and no
restriction on $\Ent(X)$.  That literal finite-$p$ reading is the only thing
tested below.  We do not claim to have established from the surrounding text
what scope the authors intended for the display, and the second line of it
carries an $o_{p\to\infty}(1)$, so an asymptotically intended reading is
consistent with everything proved here.

Two points bear on what is being refuted.  First, \eqref{eq:cand} is a hedged
candidate display: the authors write ``might be'', they do not call it a
conjecture, and we do not either.  It is nevertheless a definite mathematical
assertion, and they refer back to it, describing their first result as a
weaker form of it, so it is a statement one may test.  Second, what the
authors actually \emph{prove}, via a theorem of Tao \cite{Tao}, is the
$\varepsilon$-weakened, two-sided-window assertion
\begin{equation}\label{eq:thm}
  \Ent(X+X')\ \ge\ \Ent(X)+\tfrac12-\varepsilon
  \qquad\text{whenever } k_\varepsilon<\Ent(X)<\log p-K_\varepsilon ,
\end{equation}
a window their prose says ``forbids, e.g., deterministic random variables and
random variables uniformly distributed on the whole of $\mathbb F$''.  That
gap --- \eqref{eq:cand} has no $\varepsilon$ and no window --- is what the
witness below occupies.  Nothing in this note touches \eqref{eq:thm}; see
Section~\ref{sec:scope}.

\begin{theorem}\label{thm:main}
Let $p=101$, let $S=\{0,1,\dots,70\}\subset\mathbb F_{101}$, let $X$ be
uniform on $S$ and let $X'$ be an independent copy of $X$.  Then
\[
  \min\left\{\log p,\ \log\!\left(\frac{2^{\Ent(X)+1}-1}{\sqrt2}\right)\right\}
  \ =\ \log\!\left(\frac{141}{\sqrt2}\right)\ <\ \log 101 ,
\]
and
\[
  \Ent(X+X')\ <\ \log\!\left(\frac{141}{\sqrt2}\right).
\]
Hence \eqref{eq:cand}, read literally at this finite $p$, is false, and false
at a point where the minimum in \eqref{eq:cand} is \emph{not} the cap
$\log p$.
\end{theorem}

The last clause is the substance of the example: in the saturated branch
\eqref{eq:cand} degenerates into the demand that $X+X'$ be exactly uniform,
whereas Theorem~\ref{thm:main} places the failure strictly \emph{below} the
saturation threshold, in the branch that carries the content
$\Ent(X)+\frac12-o(1)$.

\section{The witness, exhibited}

Fix $p=101$, $m=71$, $S=\{0,\dots,m-1\}$ and $X$ uniform on $S$, so that
$\Ent(X)=\log 71=6.149747119504682\ldots$ and
\begin{equation}\label{eq:int}
  2^{\Ent(X)+1}-1\ =\ 2m-1\ =\ 141 \qquad\text{exactly.}
\end{equation}
Because $X$ is uniform on a set, the right-hand side of \eqref{eq:cand} is an
\emph{integer} expression, and this is what makes the whole test finite and
exact.

\subsection*{Which branch is active}
By \eqref{eq:int} the two branches compare as
$(2m-1)/\sqrt2$ against $p$, i.e.\ as the integers $(2m-1)^2$ against $2p^2$:
\begin{equation}\label{eq:branch}
  (2m-1)^2 = 141^2 = 19881
  \ <\ 20402 = 2\cdot 101^2 = 2p^2 .
\end{equation}
So $141/\sqrt2<101$ and the minimum in \eqref{eq:cand} is the second branch:
\begin{align*}
  \log\!\left(\frac{141}{\sqrt2}\right)&=\log 141-\tfrac12
   =6.639551352398794\ldots\\
  &<\ 6.658211482751795\ldots=\log 101 .
\end{align*}
Equivalently: the least $m$ for which the arc of length $m$ in
$\mathbb F_{101}$ saturates, i.e.\ the least $m$ with $(2m-1)^2\ge 2p^2$, is
$m=72$, and our witness has $m=71$, one step below that floor.

\subsection*{The law of $X+X'$}
Put $C_k=\#\{(i,j)\in S\times S: i+j\equiv k \bmod p\}$, so that
$\Pr[X+X'=k]=C_k/m^2$ and $\sum_k C_k=m^2$.  For the arc
$S=\{0,\dots,m-1\}$ with $p<2m-1<2p$ one has the closed form
\begin{equation}\label{eq:closed}
  C_k=f(k)+f(k+p),\qquad
  f(s)=\begin{cases}
   s+1, & 0\le s\le m-1,\\
   2m-1-s, & m-1<s\le 2m-2,\\
   0,&\text{otherwise,}
  \end{cases}
\end{equation}
since the number of ordered representations of $s$ as a sum of two elements of
$\{0,\dots,m-1\}$ \emph{over $\mathbb Z$} is $f(s)$, and each residue $k$ is
hit by exactly the integers $k$ and $k+p$ in $[0,2m-2]$.  At $p=101$, $m=71$
this gives the multiset
\begin{equation}\label{eq:multiset}
  \{C_k\}_{k=0}^{100}
  =\bigl\{\,41^{(42)},\ 42^{(2)},43^{(2)},\dots,70^{(2)},\ 71^{(1)}\,\bigr\},
\end{equation}
that is: the value $41$ occurs $42$ times, each of $42,43,\dots,70$ occurs
twice, and $71$ occurs once.  There are $42+29\cdot2+1=101$ entries and they
sum to $41\cdot42+2\cdot\sum_{v=42}^{70}v+71=1722+3248+71=5041=71^2$, as they
must.  Concretely, reading $k=0,1,\dots,100$ in order,
$C_k=41$ for $0\le k\le 40$, $C_k=k+1$ for $41\le k\le 70$, and
$C_k=141-k$ for $71\le k\le 100$.  In particular $101\,C_k\ne m^2$ for
$k=41$, so $X+X'$ is not uniform on $\mathbb F_{101}$.

For orientation only --- no step of the proof uses it ---
$\Ent(X+X')=6.6305238776\ldots$, so the deficit in \eqref{eq:cand} is
$-9.0274748\times10^{-3}$ bits.

\section{The decisive inequality, in integers}

\begin{lemma}\label{lem:crit}
Let $p$ be prime, let $S\subseteq\Fp$ with $|S|=m\ge1$, let $X$ be uniform on
$S$, let $X'$ be an independent copy and let $C_k$ be as above.  Then
\begin{equation}\label{eq:crit}
  \Ent(X+X')<\log\!\left(\frac{2m-1}{\sqrt2}\right)
  \quad\Longleftrightarrow\quad
  \bigl(2m^4\bigr)^{m^2}
  \ <\ (2m-1)^{2m^2}\Bigl(\prod_{k}C_k^{\,C_k}\Bigr)^{2},
\end{equation}
and both sides of the right-hand inequality are positive integers.
\end{lemma}

\begin{proof}
With $c_k=C_k/m^2$ and $\sum_k C_k=m^2$,
\[
  2^{\Ent(X+X')}=\prod_k c_k^{-c_k}
  =\frac{(m^2)^{\sum_k C_k/m^2}}{\bigl(\prod_k C_k^{C_k}\bigr)^{1/m^2}}
  =\frac{m^2}{P^{1/m^2}},\qquad P:=\prod_k C_k^{\,C_k},
\]
the products being over the $k$ with $C_k>0$.  Hence
$\Ent(X+X')<\log((2m-1)/\sqrt2)$ says $m^2/P^{1/m^2}<(2m-1)/\sqrt2$, i.e.
$\sqrt2\,m^2<(2m-1)P^{1/m^2}$.  Squaring (both sides are positive) gives
$2m^4<(2m-1)^2P^{2/m^2}$, and raising to the power $m^2$ gives
\eqref{eq:crit}.  Every $C_k$ is a nonnegative integer, so $P\in\mathbb Z_{>0}$.
\end{proof}

\begin{remark}\label{rem:guard}
Criterion \eqref{eq:crit} tests \eqref{eq:cand} \emph{only after} the branch
test \eqref{eq:branch} has returned ``out of band''.  In the saturated branch
the right-hand side of \eqref{eq:cand} is $\log p$, not
$\log((2m-1)/\sqrt2)$, and there the right-hand inequality of
\eqref{eq:crit} can hold with no violation whatsoever: for $S=\Fp$ itself,
$X+X'$ is uniform, \eqref{eq:cand} holds with equality, and yet
$(2m^4)^{m^2}<(2m-1)^{2m^2}P^2$ is true because it is being compared against
the \emph{inactive} branch.  So $(2m-1)^2<2p^2$ is a guard and not a
formality.  It holds for the witness of Theorem~\ref{thm:main} by
\eqref{eq:branch}.
\end{remark}

\begin{proposition}\label{prop:num}
At $p=101$, $m=71$, with $C_k$ as in \eqref{eq:multiset} and
$P=\prod_k C_k^{C_k}$, the two integers
$L=(2m^4)^{m^2}=(2\cdot71^4)^{5041}$ and
$R=(2m-1)^{2m^2}P^2=141^{10082}P^2$ satisfy
\[
  2^{91}\,L\ <\ R\ <\ 2^{92}\,L ,
\]
and in particular $L<R$.
\end{proposition}

\begin{proof}[Proof of Theorem~\ref{thm:main}]
The branch identification is \eqref{eq:int} and \eqref{eq:branch}.  The
inequality $L<R$ of Proposition~\ref{prop:num} together with
Lemma~\ref{lem:crit} gives $\Ent(X+X')<\log(141/\sqrt2)$, which is exactly
the failure of \eqref{eq:cand} at this $X$.
\end{proof}

The bracket $2^{91}L<R<2^{92}L$ is worth stating because it converts the
deficit into a rational bound with no logarithm evaluated.  Indeed the proof
of Lemma~\ref{lem:crit} shows
$\log(R/L)=2m^2\bigl(\log((2m-1)/\sqrt2)-\Ent(X+X')\bigr)$, so with
$2m^2=10082$,
\[
  \frac{91}{10082}
  \ <\ \log\!\left(\frac{141}{\sqrt2}\right)-\Ent(X+X')\ <\ \frac{92}{10082},
\]
that is, $0.0090259\ldots<\text{deficit}<0.0091251\ldots$ bits.  The violation
is therefore bounded away from zero by an inequality between integers.

\section{Further witnesses}\label{sec:table}

The same two integer tests certify a violation for the smallest out-of-band
arc at each of six more primes.  In each row $S=\{0,\dots,m-1\}$, the third
and fourth columns are the branch test $(2m-1)^2<2p^2$ of \eqref{eq:branch}, and the fifth is
and $k:=\lfloor\log_2(R/L)\rfloor$, so that $2^kL<R<2^{k+1}L$ as in
Proposition~\ref{prop:num}.  Since $k\ge0$ in every row, $L<R$ and
\eqref{eq:cand} fails at each; and by the identity displayed above the deficit
at each witness lies strictly between $k/2m^2$ and $(k+1)/2m^2$ bits.  The
last column records that bracket to three significant figures; it is the only
inexact entry in the table, and nothing depends on it.

\begin{center}
\small
\setlength{\tabcolsep}{4pt}
\begin{tabular}{rrrrrl}
$p$ & $m$ & $(2m-1)^2$ & $2p^2$ & $k$
 & deficit (bits)\\
\hline
$61$  & $43$  & $7225$   & $7442$   & $21$
 & $5.68$--$5.95\cdot10^{-3}$\\
$101$ & $71$  & $19881$  & $20402$  & $91$
 & $9.03$--$9.13\cdot10^{-3}$\\
$151$ & $106$ & $44521$  & $45602$  & $241$
 & $1.072$--$1.077\cdot10^{-2}$\\
$211$ & $147$ & $85849$  & $89042$  & $160$
 & $3.70$--$3.73\cdot10^{-3}$\\
$307$ & $213$ & $180625$ & $188498$ & $37$
 & $4.08$--$4.19\cdot10^{-4}$\\
$401$ & $278$ & $308025$ & $321602$ & $46$
 & $2.98$--$3.04\cdot10^{-4}$\\
$509$ & $353$ & $497025$ & $518162$ & $316$
 & $1.268$--$1.272\cdot10^{-3}$\\
\end{tabular}
\end{center}

\noindent
The deficit column is not monotone in $p$, and no trend should be read into
these seven rows; Section~\ref{sec:scope} returns to this.  We do not claim
that a violating arc exists at every prime.

\section{What is not settled}\label{sec:scope}

\noindent\textbf{No theorem is contradicted.}  Every witness above has
$0.50<\log p-\Ent(X)<0.53$ bits, so all seven sit at the near-uniform end of
the entropy range.  We have not determined the constant $K_\varepsilon$ of the
window of \eqref{eq:thm} and claim nothing about its size, so we do not assert
that the window excludes our witnesses; what is checked directly is that
$\Ent(X+X')-\Ent(X)>0.477$ at all seven witnesses (weakest certified bound
$0.47701$, at $p=61$), so every one of them satisfies \eqref{eq:thm} for every
$\varepsilon>0.023$.

\medskip
\noindent\textbf{Only the printed finite-$p$ form is refuted.}  What fails is
\eqref{eq:cand} exactly as displayed, with no $\varepsilon$ and no window; the
second line of the quoted display carries an $o_{p\to\infty}(1)$ and is
untouched, and a reader who takes the display to be intended asymptotically
should read this note as bearing on its finite-$p$ form alone.  We make no
claim about the $p\to\infty$ behaviour, and the seven rows of
Section~\ref{sec:table} support
none: their deficits are not monotone in $p$, and the deficit of the smallest
out-of-band arc may well tend to $0$, in which case the second line survives
intact.  Nor are the rows maximal deficits, each being the \emph{smallest}
out-of-band arc at its prime.  Whether the supremum of the deficit over all
$X$ stays bounded away from $0$ as $p\to\infty$ is open, and we do not
conjecture an answer.

\medskip
\noindent\textbf{The bulk regime is untouched, and no minimality is claimed.}
All of our witnesses sit just below the saturation floor; the regime where
$2^{\Ent(X)}$ is well below $p/\sqrt2$ is not addressed here at all.  We do
not claim that $p=101$ or $m=71$ is extremal in any sense, we search no prime
beyond $p=509$, and only arcs $\{0,\dots,m-1\}$ with uniform $X$ are
considered; no other family of sets and no non-uniform $X$ is examined.



\section{Reproduction}

Everything needed is printed above: the prime $p=101$, the set
$S=\{0,\dots,70\}$, the closed form \eqref{eq:closed} for the convolution
counts, the resulting multiset \eqref{eq:multiset}, the branch test
\eqref{eq:branch} and the integer criterion \eqref{eq:crit}.  A referee who
prefers to recompute rather than to read can do so in a few lines of exact
integer arithmetic.  A program that does exactly that, reading only the data
printed here and using the standard library alone with no floating-point
decision, accompanies this note as \texttt{verify.py}, together with the
recorded output \texttt{verify.output.txt} of a run of it.

\begin{thebibliography}{9}

\bibitem{GGK}
L.~Gavalakis, M.~Goh, and I.~Kontoyiannis,
\emph{Entropy lower bounds and sum-product phenomena},
\href{https://arxiv.org/abs/2604.20233}{arXiv:2604.20233}, version~3
(28~August 2026); quotations and the equation number refer to that version.

\bibitem{Tao}
T.~Tao,
\emph{Sumset and inverse sumset theory for Shannon entropy},
Combin. Probab. Comput. \textbf{19} (2010), no.~4, 603--639.
\href{https://doi.org/10.1017/S0963548309990642}{doi:10.1017/S0963548309990642};
\href{https://arxiv.org/abs/0906.4387}{arXiv:0906.4387}.

\end{thebibliography}

\end{document}
